%%%%%%%%%%%%%%%%%%%%%%%%%%%%%%%%%%%%%%%%%%%
\section{Vertical compactifications of adic spaces}
\label{s:compactifications}
%%%%%%%%%%%%%%%%%%%%%%%%%%%%%%%%%%%%%%%%%%%

This section is self-contained and serves as a preparation for \cref{s:fg-pi1-rig}.
Its goal is to study the universal vertical compactification $\Spa(X,S)$ of an adic space $X$ locally of finite type over a base adic space $S$, which exists under some mild hypotheses.
For rigid spaces this recovers the universal compactification from Huber's book \cite[\S 5.1]{HuberBook}, and for a map of discretely ringed adic spaces $X_0^{\rm ad}\to S_0^{\rm ad}$ induced by a map of schemes $X_0\to S_0$ one obtains the space $\Spa(X_0, S_0)$ we have already encountered in \cref{sec:adic tame discrete case}. Later in \cref{s:fg-pi1-rig}, we will need to consider a situation which is of neither kind, which is the map of adic spaces induced by an adic morphism of formal schemes $\fX\to\Spf(K^\circ)$.   The general construction, reviewed in \cref{ss:univ-comp}, is due to Solodov \cite{Solodov-vertical-compactification}.
However, in \cref{sec:root-triples} we give an alternative description of this compactification in terms of root triples. 

As we shall explain in \cref{ss:tame-pi1-rig}, in order to obtain a notion of (relative) tameness in rigid geometry with good finiteness properties, one should study finite \'etale morphisms $Y \to X$ whose extension to the universal vertical compactification $\Spa(X,S)$ satisfies a pointwise tameness condition. 

To a reader who is not interested in diving into the adic space structure of the universal vertical compactification, we recommend focusing on the definition of root triples in \cref{sec:root-triples}.
In most of the arguments in \cref{s:fg-pi1-rig} it is enough to know that the underlying set of $\Spa(X,S)$ is the set of root triples.


%%%%%%%%%%%%%%%%%%%%%%%%%%%%%%%%%%%%%%%%%%%
\subsection{Review of adic spaces}
\label{ss:adic-spaces}


We recall some general facts and terminology concerning adic spaces.
The types of adic spaces we encounter in this article are rigid spaces, their formal models, and discretely ringed adic spaces.
Accordingly, our focus will be on these spaces and their interplay.

\subsubsection*{Adic spaces}

For a Huber pair $(A, A^+)$ we denote by $A^\circ\subseteq A$ the subring of power-bounded elements and by $A^{\circ\circ}\subseteq A^\circ$ the ideal of topologically nilpotent elements. Then $(A, A^\circ)$ is a Huber pair as well, and we abbreviate $\Spa(A, A^\circ)$ to $\Spa(A)$.

If $(A, A^+)\to (B, B^+)$ is a homomorphism of Huber pairs, we write $\Spa(B, A^+)$ for the adic space $\Spa(B, B^+_{\min})$ where $B^+_{\min}\subseteq B$ is the smallest ring of integral elements containing the image of $A^+$, i.e.\ the integral closure of the sub-$A^+$-algebra of $B$ generated by $B^{\circ\circ}$.
Note that a continuous valuation of~$B$ that takes values less than or equal to~$1$ on~$A^+$ is already less than or equal to~$1$ on~$B^+_{\min}$.

An {\bf affinoid field} is a Huber pair $(K, K^+)$ where $K$ is a field and $K^+$ is a valuation ring of $K$, equipped with the valuation topology.
The property of $(K,K^+)$ being a Huber pair forces~$K^+$ to be {\bf microbial} meaning that it admits an element $\pi$ such that $K = K^+[1/\pi]$.
Every such element is called a {\bf pseudouniformiser}.
The residue field of $K^+$ is denoted by $K^\succ$ and called the {\bf specialisation field} of $K$. 

\subsubsection*{Specialisation}

Next, we discuss specialisation relations on adic spaces. A specialisation $x\leadsto x'$ on a topological space $X$ is a pair of points $x,x'\in X$ such that $x'\in\overline{\{x\}}$. 

If $x'$ is the specialisation of a point $x$ of an adic space $X$, we have a cospecialisation map $\cO_{X,x'}\to \cO_{X,x}$.
Then $x'$ is a \textbf{vertical (secondary)} specialisation of $x$ if and only if this map is a local homomorphism, so that it induces a homomorphism of residue fields $k(x')\to k(x)$. 
If $X$ is analytic, then every specialisation on $X$ is vertical. 

For a point $x\in X$ the following are equivalent:
\begin{enumerate}
    \item $x$ has no vertical generalisations,
    \item $k(x)^+ = k(x)^\circ$ (powerbounded elements),
    \item \begin{enumerate}[(i)]
            \item $x$ is analytic and $k(x)^+$ is of rank one or
            \item $k(x)$ has the trivial valuation (meaning $k(x)^+=k(x)$).
          \end{enumerate}
\end{enumerate}
We call a point satisfying the above equivalent conditions {\bf vertically maximal}.
For a point~$x$ of an adic space~$X$ we denote by~$x^\circ$ its {\bf maximal vertical generalisation}. 
It is the point corresponding to the valuation ring $k(x)^\circ$ of~$k(x)$.
Sometimes we will just call a point without vertical generalisations~$x^\circ$, without specifying any point~$x$.

\begin{lem} \label{lem:adic-preserves-vertical-maximal}
    Let $f \colon Y \to X$ be an adic morphism of adic spaces.
    Then for any point $y \in Y$ we have $f(y)^\circ = f(y^\circ)$.
    In particular,~$f$ takes vertically maximal points to vertically maximal points.
\end{lem}

\begin{proof}
    For every adic homomorphism of Huber rings $A \to B$, the image of~$A^\circ$ is contained in~$B^\circ$ (compare \cite[Lemma~1.8~(i)]{HuberContinuousValuations}).
    This implies that
    \[
        \widehat{k(f(y)^\circ)}^+ = \widehat{k(f(y))}^\circ \subseteq \widehat{k(y)}^\circ = \widehat{k(y^\circ)}^+
    \]
    and consequently, 
    \[
        \widehat{k(f(y)^\circ)}^+ \subseteq \ \widehat{k(y^\circ)}^+ \cap \widehat{k(f(y))} = \widehat{k(f(y^\circ))}^+.
    \]
    The inclusion of valuation rings translates to a vertical specialisation relation $f(y^\circ) \leadsto f(y)^\circ$.
    The maximality of~$f(y)^\circ$ (with respect to vertical generalisation) forces the equality $f(y^\circ) = f(y)^\circ$.
\end{proof}

There is a similar characterisation for horizontal specialisations: $x'$ is a {\bf horizontal (primary) specialisation} if and only if $\cO_{X,x'}^+\to\cO^+_{X, x}$ is a local homomorphism, so that we get a homomorphism of specialisation fields $k(x')^\succ \to k(x)^\succ$.
Finally, a {\bf generalised horizontal specialisation} is a specialisation $x \leadsto x'$ that is either horizontal or factors as $x \leadsto y \leadsto x'$ such that~$y$ and~$x'$ are trivial valuations, and $x \leadsto y$ is horizontal.


\subsubsection*{Rigid spaces}

A {\bf non-archimedean field} is a field $K$ which is complete with respect to a rank one valuation. Its valuation ring $K^+ = K^\circ$ is thus $\pi$-adically complete with respect to every pseudouniformiser $\pi$ of $K^\circ$.

Let $K$ be a non-archimedean field. A {\bf rigid space} over $K$ is an adic space locally of finite type over $\Spa(K)$. So, a rigid space over $K$ is locally of the form $\Spa(A)$ for an affinoid $K$-algebra $A$. We shall denote the category of rigid spaces over $K$ by $\cat{Rig}_K$.


\subsubsection*{Formal schemes}

For our subsequent discussion of formal models and the specialisation map we shall rely on \cite{FujiwaraKato}, using \cite[Theorem~II~A.5.1]{FujiwaraKato} to translate results to the language of adic spaces.

Let $K$ be a non-archimedean field with valuation ring $K^+=K^\circ$ and pseudouniformiser $\pi$. If $A$ is an algebra topologically of finite type over $K^\circ$ (that is, a quotient of $K^\circ\langle T_1, \ldots, T_n\rangle$ for some $n\geq 1$, endowed with the $\pi$-adic topology), then $A_K = A[1/\pi]$ is an affinoid $K$-algebra and $A_K^\circ$ is the integral closure of the (image of) $A$ in $A_K$. We call a formal scheme over $K^\circ$ locally of finite type if it is locally of the form $\Spf(A)$ for a $K^\circ$-algebra $A$ topologically of finite type. We denote the category of such formal schemes over $K^\circ$ by $\cat{FSch}^{\rm lft}_{K^\circ}$. 

In \cite[\S~4]{Huber94:Generalization} Huber explains how the category of locally noetherian formal schemes embeds fully faithfully into the category of adic spaces by sending $\Spf(A)$ to $\Spa(A)$.
Being noetherian was needed because he had established sheafiness of~$A$ only in the noetherian setting (see \cite[Theorem~2.2]{Huber94:Generalization}).
Today we have \cite{ZavyalovSheafy} at hand that settles sheafiness also if~$A$ is only strongly rigid noetherian (e.g. topologically of finite type over a microbial valuation ring).
Consequently, we have a fully faithful functor
\[
    (-)^{\rm ad} \colon \cat{FSch}^{\rm lft}_{K^\circ} \longrightarrow \cat{AdSpc}_{K^\circ}
\]
to the category of adic spaces that are locally of finite type over~$K^\circ$.

\subsubsection*{Rigid generic fibre}

There exists a unique functor 
\[
    (-)_{\rm rig} \colon \cat{FSch}^{\rm lft}_{K^\circ} \longrightarrow \cat{Rig}_K,
\]
called the {\bf rigid generic fibre} functor, which sends open coverings to open coverings, and such that 
\[ 
    \Spf(A)_{\rm rig} \simeq \Spa(A_K).
\]
If~$\cX$ is a formal scheme locally of finite type over~$K^\circ$, then~$\cX_{\rm rig}$ is the generic fibre of the morphism $\cX^{\rm ad} \to \Spa (K^\circ)$.
Being a fibre product of adic spaces, it preserves open and closed immersions, finite maps, étale maps, etc.

\subsubsection*{The specialisation map}

Let $A$ be an algebra topologically of finite type over $K^\circ$. For any locally topologically ringed space $X$, we have a functorial isomorphism $\Hom(X, \Spf(A))\simeq \Hom(A, \cO(X))$ (continuous homomorphisms). Thus, the map $A\to A_K^\circ$ induces a map of locally topologically ringed spaces
\[ 
    {\rm sp}_A \colon (\Spf(A)_{\rm rig}, \cO^+) = (\Spa(A_K),  \cO^+) \longrightarrow (\Spf(A), \cO). 
\]
This construction respects basic open subsets: the preimage of $\{f\neq 0\}=\Spf(A\langle f^{-1}\rangle)$ is $\{|f|\geq  1\} = \Spa(A\langle f^{-1}\rangle_K$ and ${\rm sp}_A$ restricts to ${\rm sp}_{A\langle f^{-1}\rangle}$.  Hence it globalises: for every formal scheme $\fX$ locally of finite type over $K^\circ$ there exists a unique map of topologically locally ringed spaces
\[
    {\rm sp}_{\fX} \colon (\fX_{\rm rig}, \cO^+)\longrightarrow \fX
\]
such that for every map $\Spf(A)\to \fX$ of formal schemes locally of finite type over $K^\circ$, the resulting square 
\[ 
    \begin{tikzcd}
        (\Spa(A_K), \cO^+) \ar[d] \ar[r,"{\rm sp}_A"] & \Spf(A) \ar[d] \\
        (\fX_{\rm rig}, \cO^+) \ar[r,"{\rm sp}_\fX",swap] & \fX.
    \end{tikzcd}
\]
commutes.
Furthermore, we have ${\rm sp}_{\fX,*}(\cO) = \cO_\fX[1/\pi]$, and ${\rm sp}_{\fX,*}(\cO^+)$ coincides with the integral closure of (the image of) $\cO_{\fX}$ in $\cO_{\fX}[1/\pi]$.

\begin{defi}[{\cite[Definition~A.1]{AchingerLaraYoucis2022}}] \label{defi:eta-normal}
    A formal scheme $\fX$ locally of finite type over~$K^\circ$ is 
    {\bf \mbox{$\eta$-normal}} if $\cO_{\fX}$ is $\pi$-torsion free and integrally closed in $\cO_{\fX}[1/\pi]$, or equivalently if the map $\cO_\fX \isomto {\rm sp}_{\fX*}(\cO^+)$ is an isomorphism.
\end{defi}

\begin{rmk} \label{rmk:sp-facts}
We shall use the following facts about ${\rm sp}={\rm sp}_\fX\colon X=\fX_{\rm rig}\to \fX$.
\begin{enumerate}
    \item On the level of points, the specialisation map can be described as follows. Let $x\in X$ and let $\Spf(A)\subseteq\fX$ be an affine neighbourhood of ${\rm sp}(x)$, so that $\Spa(A_K)\subseteq X$ is an affinoid neighbourhood of $x$. Then $k(x)^+$ is a $\pi$-adically separated valuation ring, and the map $A\to k(x)^+$ induces a map of formal schemes 
    \[
        \Spf(k(x)^+) \longrightarrow \Spf(A) \subseteq \fX
    \]
    which sends the closed point of $\Spf(k(x)^+)\simeq \Spec(k(x)^+/\sqrt{(\pi)})$ to ${\rm sp}(x)$.
    \item If $\fX$ is {\bf admissible} (i.e. flat over $K^\circ$), then ${\rm sp}$ is surjective \cite[Proposition II 3.1.5]{FujiwaraKato}. 
    \item If $\fX$ is moreover $\eta$-normal, then every generic point $x\in \fX$ has a unique preimage $x'\in X$ under ${\rm sp}$. Moreover, we have
    \[
        \cO_{\fX, x} = \cO^+_{X, x'} = k(x')^+,
    \]
    and hence $k(x) = k(x')^\succ$. In particular, $\cO_{\fX,x}$ is a rank one valuation ring. These assertions follow by combining \cite[Proposition~2.2]{BhattHansen} with \cite[Lemma~A.12 and Proposition~A.15]{AchingerLaraYoucis2022}. 
\end{enumerate}
\end{rmk}

\subsubsection*{Raynaud's equivalence}

The rigid generic fibre functor sends the class $W$ of {\bf admissible formal blow-ups} (completed blowups in open coherent ideals) to isomorphisms, and induces by \cite[Theorems~II~A.5.1,~A.5.2]{FujiwaraKato} an equivalence
\[ 
    \cat{FSch}^{\rm ft}_{K^\circ}[W^{-1}] \isomto \cat{Rig}_K^{\rm qcqs}
\]
between formal schemes of finite type over $K^\circ$ with admissible formal blow-ups formally inverted and quasi-compact and quasi-separated rigid spaces over $K$. Both categories are equivalent to $\cat{FSch}^{\rm adm}_{K^\circ}[W^{-1}]$ where $\cat{FSch}^{\rm adm}_{K^\circ}$ is the category of admissible formal schemes over $K^\circ$. 

\subsubsection*{Formal models}

Let $X$ be a rigid space over $K$. A {\bf formal model} (resp.\ {\bf admissible formal model}) of $X$ is a formal scheme $\fX$ locally of finite type over $K^\circ$ (resp.\ admissible) together with an identification $\fX_{\rm rig}\simeq X$. We denote by $M(X)$ (resp.\ $M_{\rm adm}(X)$) the category of formal models (resp.\ admissible formal models) of $X$. Then $M(X)$ is a cofiltering category and $M_{\rm adm}(X)$ is (equivalent to) a cofiltering poset. The specialisation maps ${\rm sp}_\fX\colon X\to \fX$ for $\fX\in M(X)$ commute with the maps in $M(X)$ and induce an isomorphism of locally ringed spaces
\[
    (X, \cO_X^+)\isomto \varprojlim_{\fX\in M(X)} \, \fX. 
\]
In particular, for every $x\in X$ we have $\cO^+_{X,x} = \varinjlim_{\fX\in M(X)} \cO_{\fX, {\rm sp}_{\fX}(x)}$.
In fact, the admissible formal models $M_{\rm adm}(X)$ are cofinal in~$M(X)$ so it suffices to take the direct limit over~$M_{\rm adm}(X)$.
Moreover, if we start with an admissible formal model~$\cX$, we obtain a cofinal class of formal models of~$X$ by taking all admissible formal blow-ups of~$\cX$, see \cite[Chapter~II,\S~3.2 and A.5]{FujiwaraKato}.

%%%%%%%%%%%%%%%%%%%%%%%%%%%%%%%%%%%%%%%%%%%
\subsection{Vertical compactifications of adic spaces}
\label{ss:univ-comp}

In \cite[\S5.1]{HuberBook} Huber constructs the universal compactification for a given morphism $f\colon Y \to X$ of analytic adic spaces subject to the following hypotheses.
According to his general assumptions, the spaces~$X$ and~$Y$ have to be locally noetherian and, moreover, the morphism~$f$ is required to be locally of ${}^+$weakly finite type\footnote{Locally of the form $\Spa(B,B^+)\to\Spa(A,A^+)$ where $A\to B$ is topologically of finite type and where $B^+$ is the smallest integral subring containing a finitely generated $A^+$-algebra, see \cite[Definition 1.2.1]{HuberBook}.} , separated, and taut\footnote{A locally spectral topological space is {\bf taut} if it is quasi-separated and every quasi-compact open has quasi-compact closure. A spectral map $f\colon Y\to X$ is {\bf taut} if the preimage of every taut open of $X$ is taut. Examples include qcqs and partially proper morphisms of adic spaces. See \cite[\S 5.1]{HuberBook}.}.
The situation is quite different from the scheme setting, where we normally do not have an initial object among all compactifications.
What makes analytic adic spaces special in this respect is the fact that all specialisations are vertical.
The universal compactification is constructed by adding all possible vertical specialisations; there is no choice involved.

If we want to drop the hypothesis that~$X$ and~$Y$ are analytic, we have to account for horizontal specialisations entering the picture.
As a consequence, we can no longer expect the existence of a universal compactification.
However, we can construct the \emph{universal vertical compactification} as in \cite{Solodov-vertical-compactification} that can be thought of as adding only \emph{vertical} specialisations.
In this way, we obtain a \emph{vertically partially proper} morphism, which is a morphism that satisfies a valuative criterion for properness only for vertical specialisations.
Let us start by explaining this notion.

Recall from \cite[Definitions~8.26 and~8.54]{Wedhorn-adicspaces} that a Huber ring~$A$ is {\bf stably sheafy} if every Huber pair $(B,B^+)$ such that $B$ is topologically of finite type over~$A$ is sheafy.
An adic space is {\bf stable} if it is locally of the form $\Spa(A,A^+)$ for a stably sheafy Huber ring~$A$.

\begin{defi}
    Let~$S$ be a stable adic space and $f \colon X \to S$ a morphism locally of weakly finite type. 
    \begin{enumerate}[(i)]
        \item The morphism~$f$ is {\bf vertically specialising} at a point $x \in X$ if every vertical specialisation $f(x) \rightsquigarrow s'$ in~$S$ can be lifted to a specialisation $x \rightsquigarrow x'$ in~$X$.
        If~$f$ is vertically specialising at every point, we call it {\bf vertically specialising}.
        \item We call~$f$ {\bf universally vertically specialising} if every base change of~$f$ along an adic  morphism from a stable adic space is vertically specialising.
    \end{enumerate}
\end{defi}

In the above definition the lifted specialisation $x \rightsquigarrow x'$ is not required to be vertical.
However, if a lift exists, we can always modify it to a \emph{vertical} specialisation by \cite[Proposition~3.4]{Solodov-vertical-compactification}.
We further remark that in the definition of universally vertically specialising morphisms we test after base change via an \emph{adic} morphism because then we can guarantee that the respective fibre product exists.

\begin{lem}[Valuative criterion for being universally vertically specialising] \label{valuative-criterion-separatednes}
    Let~$S$ be a stable adic space and $f \colon  X \to S$ a morphism locally of weakly finite type.
    Then~$f$ is universally vertically specialising if and only if for every solid arrow diagram 
    \begin{equation} \label{eqn:valuative-criterion-test-square} 
        \begin{tikzcd}
            \Spa(k, k_2^+) \ar[r] \ar[d] & X \ar[d] \\
            \Spa(k, k_1^+) \ar[r] \ar[ur,dashed] & S,
        \end{tikzcd}
    \end{equation}
    where $(k,k_1^+)$ and~$(k,k_2^+)$ are affinoid fields with $k_1^+ \subseteq k_2^+$, there is a diagonal dashed arrow making the diagram commutative.
\end{lem}

\begin{proof}
    Analogous to \cite[Proposition~1.3.8]{HuberBook}, see also \cite[Remark~3.5]{Solodov-vertical-compactification}. 
\end{proof}

For simplicity we define vertical separatedness directly in terms of a valuative criterion.

\begin{defi}
    Let~$S$ be a stable adic space and $f \colon  X \to S$ a morphism that is locally of weakly finite type.
    Then~$f$ is {\bf vertically separated} if it is quasi-separated and for every solid arrow square (\ref{eqn:valuative-criterion-test-square}) there is at most one dashed arrow making the diagram commutative.
\end{defi}

As is to be expected, separatedness implies vertical separatedness (see the argument in \cite[Proposition~1.3.8]{HuberBook}) and a universally closed morphism is universally specialising and thus universally vertically specialising.

\begin{defi}
    Let~$S$ be a stable adic space.
    A morphism $f \colon  X \to S$ is {\bf vertically partially proper} if it is locally of ${}^+$weakly finite type, vertically separated and universally vertically specialising.

    It is {\bf vertically proper} if it is vertically partially proper and quasi-compact.
\end{defi}

Let~$S$ be a stable adic space.
It follows from \cref{valuative-criterion-separatednes} and the definition of vertical separatedness that a quasi-separated morphism $f \colon X \to S$ that is locally of ${}^+$weakly finite type is vertically partially proper if and only if for every solid arrow square (\ref{eqn:valuative-criterion-test-square}) there is a unique dashed arrow.

\begin{defi}
    Let~$S$ be a stable adic space and $f \colon  X \to S$ a morphism locally of ${}^+$weakly finite type.
    A {\bf vertical compactification} of~$f$ is a commutative triangle
    \[
     \begin{tikzcd}
         X  \ar[rr,open,"j"] \ar[dr,"f"']   &       & \ov{X}    \ar[dl,"\bar{f}"]  \\
                                            & S
     \end{tikzcd}
    \]
    of adic spaces in which $j$ is a locally closed embedding and~$\bar{f}$ is vertically partially proper.
    We call~$\overline{X}$ the {\bf universal vertical compactification} if, in addition, it is initial among all morphisms over~$S$ to adic spaces that are vertically partially proper over~$S$.
\end{defi}

\begin{rmk}
    In \cite{Solodov-vertical-compactification}, Solodov requires vertically partially proper morphisms to be separated (and not just vertically separated).
    Accordingly, at first sight, his universal vertical compactifications satisfy a slightly weaker universal property that only takes test objects $Y \to S$ that are vertically partially proper \emph{and} separated.
    However, the universal vertical compactification he constructs also satisfies the universal property without the separatedness assumption.
    Below we will retrace his arguments and explain why this is the case.
\end{rmk}

Generalising \cite[Lemma~5.1.7]{HuberBook}, Solodov gives sufficient conditions for a morphism $\ov X \to S$ to be the universal compactification of $X \to S$.
Here, we give a slightly stronger version because in our terminology vertical partial properness does not include separatedness (only vertical separatedness).

\begin{prop}[\cite{Solodov-vertical-compactification}, Lemma~3.19~iv)] \label{prop:check-univprop-comp}
    We consider a commutative triangle of adic spaces
    \[
        \begin{tikzcd}
            X \ar[r,"j",open] \ar[d,"f"]    & \ov X \ar[dl,"\bar{f}"]   \\
            S
        \end{tikzcd}
    \]
    with the following properties:
    \begin{itemize}
        \item $\bar{f}$ is vertically partially proper
        \item $j$ is a quasi-compact open immersion,
        \item $X$ is closed under generalised horizontal specialisations in~$\ov X$,
        \item every point of~$\ov X$ is a vertical specialisation of a point of~$X$,
        \item $\cO_{\ov X} \to j_*\cO_X$ is an isomorphism of sheaves of topological rings.
    \end{itemize}
    Then the above commutative triangle exhibits~$\ov X$ as the universal vertical compactification of~$X$ over~$S$ and this assertion even holds after base change to some open subspace of~$S$.
\end{prop}

\begin{proof}
    Solodov's proof in \cite[Lemma~3.19~iv)]{Solodov-vertical-compactification} basically just combines items~i) and~iii) of his Lemma~3.19, which only require vertical separatedness.
\end{proof}

In \cref{sec:root-triples} we will identify the universal vertical compactification~$\overline{X}$ of $f \colon X \to S$ (if it exists, and under mild supplementary geometric assumptions) with an explicitly defined space $\Spa(X,S)$. 
But before having established that $\Spa(X,S)$ is indeed the universal vertical compactification of~$f$, we will exclusively use~$\overline{X}$ to denote the universal vertical compactification.

In order to ensure the existence of the universal vertical compactification of $f \colon X \to S$, some mild additional assumptions about~$S$ and~$X$ are necessary.
The condition is called \emph{(weak) square completeness}, a concept introduced in \cite[Definition~3.13]{Huebner2025:ProperBaseChange} and further differentiated in \cite[Definition~3.20]{Solodov-vertical-compactification}.
\begin{defi}
    An adic space~$X$ is called {\bf weakly square complete} if for three given points $x$, $x_h$, and $x_v$ in~$X$ such that~$x_h$ is a horizontal and~$x_v$ a vertical specialisation of~$x$, there is a point $x' \in X$ that is a horizontal specialisation of~$x_v$ and a vertical specialisation of~$x_h$.
     \[
      \begin{tikzcd}
       x_v	\ar[r,rightsquigarrow]							& x'			\\
       x	\ar[r,rightsquigarrow]	\ar[u,rightsquigarrow]	& x_h	\ar[u,rightsquigarrow]
      \end{tikzcd}
     \]
     If the point~$x'$ is unique for any choice of~$x$, $x_h$, and~$x_v$, we call~$X$ {\bf square complete}.
\end{defi}

Examples of square complete adic spaces include affinoid spaces, analytic spaces, formal schemes, and spaces of the form $\Spa(X,S)$ for a morphism of schemes $X \to S$. For an example of a non \emph{weakly} square-complete discretely ringed adic space see \cite[Example~3.12]{Huebner2025:ProperBaseChange}. 
We are now ready to state the existence theorem of the universal vertical compactification.
Its construction has been worked out in \cite{Solodov-vertical-compactification}.

\begin{thm}[{\cite[Theorem~3.29]{Solodov-vertical-compactification}}] \label{thm:Solodov}
    Let~$S$ be a square complete and stable adic space and $f\colon  X \to S$ a separated and taut morphism that is locally of ${}^+$weakly finite type and such that~$X$ is weakly square complete.
    Then~$f$ admits a universal vertical compactification.
\end{thm}

The author, Ronald Solodov, closely follows Huber's line of argumentation for the construction of the universal compactification of an analytic adic space.
Note that even in the analytic setting, Solodov's result is more general than Huber's as it does not make any noetherianness assumptions.

\smallskip

In the remaining part of this section we will make the construction of the universal compactification more explicit.
Let us start with the case of a morphism of affinoid adic spaces.

\begin{prop}[{\cite[Theorem~3.14]{Solodov-vertical-compactification}}] \label{prop:affinoid-compactification}
    Let $f\colon X\to S$ be a morphism of ${}^+$-weakly finite type between affinoid adic spaces $X = \Spa(A,A^+)$ and $S = \Spa(R,R^+)$ where $S$ is stable.
    Then the universal vertical compactification of~$f$ is given by
    \[
     \bar{f} \colon \overline{X} = \Spa(A,R^+) \longrightarrow \Spa(R,R^+) = S.
    \]   
\end{prop}

In the general case, Solodov (following Huber) first covers~$S$ by affinoid open subspaces~$S_i$ and shows that universal vertical compactifications over~$S_i$ can be glued together to a universal vertical compactification over~$S$.
This uses the last part of \cref{prop:check-univprop-comp}.

Consequently, it suffices to construct~$\ov X$ in case~$S$ is affinoid.
Now Solodov covers~$X$ with affinoid subspaces~$X_i$ and aims at gluing the respective universal compactifications~$\overline{X}_i$.
Here, a major technical difficulty arises because we cannot expect the spaces~$\overline{X}_i$ to be open in the universal compactification~$\overline{X}$ of~$X$.
Huber overcomes these difficulties by finding appropriate open subspaces~$U_i$ of~$\overline{X}_i$ that are expected to be open in~$\overline{X}$ but are large enough to cover~$\overline{X}$.
The universal compactification~$\overline{X}$ is then constructed by gluing the subspaces~$U_i$.
Solodov generalises this approach for his construction of the universal vertical compactification.
As a result, we are guaranteed the existence of the universal vertical compactification.
In addition, we extract from its proof the existence of the open subsets~$U_i$ mentioned above:

\begin{lem} \label{lem:interiors-cover}
    In the setting of \cref{thm:Solodov}, assume that~$S$ is affinoid. 
    \begin{enumerate}[(1)]
        \item \label{lemitem:interior-open}
        Let $U \subseteq X$ be an open affinoid subspace and consider the injective map $\ov{U} \to \ov{X}$ of universal vertical compactifications. Let $\ov{U}^\circ$ be the subset of $\ov{U}$ consisting of all points $y \in \ov{U}$ such that for any vertical generalisation~$x$ of~$y$ in $\ov{X}$ we either have $x \in U$ or $x \notin X$. 
        Then $\ov{U}^\circ$ is an open subset of $\ov{X}$ such that $\ov{U}^\circ \cap X = U$.

        \item \label{lemitem:interior-open-covering}
        Let $X = \bigcup_i U_i$ be an open covering by affinoid subspaces. Then the subsets $\ov{U_i}^\circ$ as constructed in \labelcref{lemitem:interior-open} associated to $U_i$ form an open covering of $\ov{X}$.
    \end{enumerate}
\end{lem}

\begin{proof}
    \labelcref{lemitem:interior-open} In \cite{Solodov-vertical-compactification}, Solodov considers the universal vertical compactification~$\ov U$ (there denoted by~$U_c$) of an affinoid open $U \subseteq X$.
    He constructs a certain subspace $U_d \subseteq \ov U$ containing~$U$ that turns out to be open.
    In addition, an inclusion of open affinoids $U \subseteq V \subseteq X$ induces an open immersion $U_d \to V_d$, see \cite[Lemma~3.27]{Solodov-vertical-compactification}.
    Looking at the definition of~$U_d$ we identify it with $\ov{U}^\circ$ as the subset of~$\ov U$. The identification $\ov{U}^\circ \cap X = U$ follows from the definition of~$\ov{U}^\circ$.

    \labelcref{lemitem:interior-open-covering}
    In the end~$\ov X$ is obtained by gluing the spaces~$\ov{U_i}^\circ = U_{i,d}$ (in Solodov's notation) constructed from the members~$U_i$ of the open covering.
    Then the spaces~$\ov{U_i}^\circ$ form an open covering of~$\ov X$ by construction.   
\end{proof}


In the end, Solodov checks that his space~$\ov X$ satisfies the universal property using \cite[Lemma~3.19~iv)]{Solodov-vertical-compactification}.
With the updated version \cref{prop:check-univprop-comp} of this criterion we ensure that we can use any vertically partially proper morphism as test object (it does not need to be separated).

%%%%%%%%%%%%%%%%%%%%%%%%%%%%%%%%%%%%%%%%%%%%%%%%
\subsection{Root triples} 
\label{sec:root-triples}

Let us sketch an alternative description of the universal vertical compactification that allows for an explicit definition of the points of $\overline{X}$.
In what follows, we will write down a set $\Spa(X,S)$ and endow it with a topology and a structure sheaf.
This construction works without any hypotheses on~$X$ or~$S$.
Our final goal will be to identify this space $\Spa(X,S)$ under the assumptions of \cref{thm:Solodov} with the universal vertical compactification~$\overline{X}$, in this way also showing that it carries the structure of an adic space. 

\begin{defi}
    Let $f\colon X\to S$ be a morphism of adic spaces. A \textbf{point of $X$ with centre on $S$} is a commutative square of the form
    \begin{equation} \label{eqn:pp-test-square} 
        \begin{tikzcd}
            \Spa(k, k^\circ) \ar[r,"u"] \ar[d] & X \ar[d] \\
            \Spa(k, k^+) \ar[r,"v",swap] & S
        \end{tikzcd}
    \end{equation}
    where $(k, k^+)$ is an affinoid field with ring of powerbounded elements $k^\circ$ (so $(k, k^\circ)$ is an affinoid field as well), and where the left vertical map is induced by the inclusion $k^+\subseteq k^\circ$ and~$u$ is adic. We denote such a datum by $(k^+, u, v)$ for simplicity. 

    We deem two points $(k_i^+, u_i, v_i)$, for $i=1,2$, of $X$ with centre on $S$ \textbf{equivalent} if there exists a third square with $(k^+, u, v)$ and commutative diagrams
    \[ 
        \begin{tikzcd}
            \Spa(k_i, k^\circ_i) \ar[r] \ar[d] \ar[rr,"u_i",bend left=15] & \Spa(k, k^\circ) \ar[r,"u",swap] \ar[d] & X \ar[d]  \\
            \Spa(k_i, k^+_i) \ar[r] \ar[rr,"v_i",swap,bend right=15] & \Spa(k, k^+)\ar[r,"v"] & S
        \end{tikzcd}
    \]
    for $i=1,2$ such that $(k_i,k_i^+)/(k,k^+)$ are extensions of valuation rings, i.e.\ $k^+=k_i^+\cap k$.
\end{defi}

By \cite[Lemma~4.3]{Solodov-vertical-compactification} this defines an equivalence relation.
In every equivalence class we can identify a distinguished representative given by a so called root triple.
Before we can define root triples, we need to recall centres of valuations.
For a pair $(x, V)$ consisting of a point $x\in X$ and a valuation subring $V\subseteq k(x)^+$, by a {\bf centre of $(x, V)$ on $S$} we shall mean a vertical specialisation $s\in S$ of $f(x)$ with the property that $k(s)^+ = V\cap k(s)$.
The centres of $(x,V)$ on~$S$ are in one to one correspondence with the morphisms $\Spa(k(x),V) \to S$ making the diagram
\[
 \begin{tikzcd}
     \Spa(k(x),k(x)^+)  \ar[r]  \ar[d]  & X \ar[d]  \\
     \Spa(k(x),V)       \ar[r,dotted]          & S
 \end{tikzcd}
\]
commutative.

\begin{defi} \label{defi:root-triple}
    A {\bf root triple} of $X/S$ is a triple $x = (x^\circ,k(x)^+,x^\succ)$ for a point $x^\circ \in X$ without vertical generalisations and a valuation subring 
    \[
        k(x)^+ \subseteq k(x^\circ)^+ = k(x^\circ)^\circ
    \]
    with centre $x^\succ \in S$. For a root triple~$x$ as above we use the notation $k(x)$ to denote $k(x^\circ)$.
    Moreover, we sometimes write $x^\succ$ also for the induced map $\Spa(k(x),k(x)^+) \to S$.
\end{defi}

A root triple $x = (x^\circ,k(x)^+,x^\succ)$ naturally defines a point of~$X$ with centre on~$S$
    \[  
        \begin{tikzcd}
            \Spa(k(x), k(x)^\circ) \ar[r] \ar[d] & X \ar[d] \\
            \Spa(k(x), k(x)^+) \ar[r,"x^\succ"] & S
        \end{tikzcd}
    \]
and indeed every equivalence class of points of~$X$ with centre on~$S$ has a unique representative given by a root triple (see \cite[Lemma~4.2]{Solodov-vertical-compactification}).
In particular, this ensures that the equivalence classes form a set.

\begin{defi}
    We denote the set of equivalence classes of points of~$X$ with centre on~$S$ by $\Spa(X,S)$.
\end{defi}

In what follows, we shall often implicitly identify a point $x \in \Spa(X,S)$, which is a priori an equivalence class of squares, with the corresponding root triple and we write $x = (x^\circ,k(x)^+,x^\succ)$.

\begin{ex} \label{ex:comp-affinoid}
    For the identity morphism $X \to X$ we obtain a natural identification of sets
    \[
     X \overset{\sim}{\longrightarrow} \Spa(X,X), \qquad x \mapsto (x^\circ,k(x)^+,x).
    \]
    The inverse sends a root triple $(x^\circ,k(x)^+,x^\succ)$ to~$x^\succ$.
\end{ex}

The formation of $\Spa(X,S)$ enjoys the following functoriality.
If
    \[ 
        \begin{tikzcd}
            X'\ar[d] \ar[r,"f_X"] & X\ar[d] \\
            S'\ar[r,"f_S"] & S
        \end{tikzcd}
    \]
is a commutative square of adic spaces such that~$f_X$ is adic, we obtain a natural morphism
\[
 \Spa(f_X,f_S) \colon \Spa(X',S') \longrightarrow \Spa(X,S)
\]
by sending a point $(k^+,u,v)$ of~$X$ with centre on~$S$ to $(k^+,f_X \circ u,f_S \circ v)$.
\begin{equation} \label{diagam:functoriality-comp}
 \begin{tikzcd}
     \Spa(k,k^\circ)    \ar[r,"u"]  \ar[d]  & X'    \ar[r,"f_X"]    \ar[d]  & X \ar[d]  \\
     \Spa(k,k^+)        \ar[r,"v"]          & S'    \ar[r,"f_S"]            & S.
 \end{tikzcd}
\end{equation}
This construction is compatible with the equivalence relation, and thus indeed defines a map $\Spa(X',S') \to \Spa(X,S)$.
In terms of root triples it sends $x' = (x'^\circ,k(x')^+,x'^\succ)$ to
\[
    f(x') = (f_X(x'^\circ),k(x')^+ \cap k(f(x'^\circ)),f_S(x'^\succ)).
\]
Note that this implicitly uses \cref{lem:adic-preserves-vertical-maximal} to ensure that $f(x'^\circ)$ is vertically maximal.

If~$f_X$ is not adic, we still get a map $\Spa(f_X,f_S)$ as above but its description is a bit more complicated.
Starting with a point $(k^+,u,v)$ we can consider the diagram~(\ref{diagam:functoriality-comp}) but $f_X \circ u$ might not be adic, in which case it does not directly define a point of~$X$ with centre on~$S$.
The morphism $f_X \circ u$ can only possibly not be adic when $(k,k^\circ)$ is analytic but the image point~$x$ of $f_X \circ u$ is not analytic.
In this case we replace the diagram with
\[
    \begin{tikzcd}
        \Spa(k,k)   \ar[r,"u"]  \ar[d]  & X \ar[d]  \\
        \Spa(k,k^+) \ar[r,"v"]          & S,
    \end{tikzcd}
\]
where the affinoid fields in the left column are now endowed with the discrete topology.
Since~$x$ is nonanalytic, the morphisms are still well defined and the above diagram defines a point of~$X$ with centre on~$S$.
We declare this point to be the image of $(k^+,u,v)$ in $\Spa(X,S)$.

The root triple description is obtained as follows:
We first note that we always have $k(f(x'^\circ)^\circ) = k(f(x'^\circ))$ because $f(x'^\circ)^\circ \ne f(x'^\circ)$ can only happen if these points are nonanalytic, in which case the residue fields are discrete.
We send $x' = (x'^\circ,k(x')^+,x'^\succ)$ to the root triple $x$ where $x^\circ = f(x'^\circ)^\circ$, $k(x)^+ = k(x')^+ \cap k(f(x'^\circ))$, and $x^\succ = f_S(x'^\succ)$.

\begin{lem} \label{lem:Spa-injective}
    If $f_X \colon X' \to X$ is a locally closed embedding and $f_S \colon S' \to S$ is separated, the map $\Spa(f_X,f_S)$ is injective.
\end{lem}

\begin{proof}
    Suppose that $x_1$ and $x_2$ are two points of $\Spa(X',S')$ that map to the same point in $\Spa(X,S)$.
   We identify~$x_i$ with the corresponding root triple $x_i = (x_i^\circ,k(x_i)^+,x_i^\succ)$ and consider the diagrams
    \[
 \begin{tikzcd}
     \Spa(k(x_i),k(x_i)^\circ)    \ar[r]        \ar[d]  & X'    \ar[r,"f_X"]    \ar[d]  & X \ar[d]  \\
     \Spa(k(x_i),k(x_i)^+)        \ar[r,"x_i^\succ"]          & S'    \ar[r,"f_S"]            & S.
 \end{tikzcd}
\]
Since~$f_X$ is injective, we obtain $x_1^\circ = x_2^\circ$.
In particular, the upper line is the same in both diagrams.
We denote by~$x^\circ$ the image in~$X$ of $x_1^\circ = x_2^\circ \in X'$.
The morphism~$f_X$ being a locally closed embedding, the residue field extension $k(x_i)/k(x)$ induces an isomorphism on completions.
In particular, $k(x_i)^+$ is uniquely determined by $k(x_i)^+ \cap k(x)$, which is the same for $i=1,2$.
Finally, the morphisms~$x_1^\succ$ and~$x_2^\succ$ coincide by the valuative criterion for separatedness applied to the square
\[
    \begin{tikzcd}
        \Spa(k(x_1),k(x_1)^\circ) = \Spa(k(x_2),k(x_2)^\circ)   \ar[r]  \ar[d]  & S' \ar[d,"f_S"]    \\
        \Spa(k(x_1),k(x_1)^+) = \Spa(k(x_2),k(x_2)^+)           \ar[r]  & S,
    \end{tikzcd}
\]
where the lower map is $f_S \circ x_1^\succ = f_S \circ x_2^\succ$.
\end{proof}

As a next step we want to equip the set $\Spa(X,S)$ of equivalence classes of points of~$X$ with centre on~$S$ with a topology.
In the affinoid case this is straightforward.

\begin{lem} \label{lem:Spa-barX-affinoid}
    If $X = \Spa(A,A^+)$ and $S = \Spa(R,R^+)$ are affinoid, there is a natural identification (as sets)
    \[
    \Spa(X,S) = \Spa(A,R^+)
    \]
    and in case $S$ is stable and $X \to S$ is ${}^+$weakly of finite type, $\Spa(A,R^+)$ is the universal vertical compactification of~$X$ over~$S$.
\end{lem}

\begin{proof}
    An element of $\Spa(X,S)$ is defined by a root triple $x = (x^\circ,k(x)^+,x^\succ)$.
    Giving the pair $(x^\circ,k(x)^+)$ is equivalent to giving a continuous valuation $v$ on~$A$.
    The centre~$x^\succ$ is unique in the affinoid case.
    Its existence is equivalent to the valuation~$v$ taking values less than or equal to~$1$ on~$R^+$.
    The last assertion of the lemma follows from  \cite[Theorem~3.14]{Solodov-vertical-compactification}.
\end{proof}

By \cref{lem:Spa-barX-affinoid} the set $\Spa(X,S)$ for affinoid spaces~$X$ and~$S$ naturally carries the structure of an adic space.
In particular, it is endowed with a topology.
It is clear from the construction that commutative squares
\[
 \begin{tikzcd}
     X' \ar[r]  \ar[d]  & X \ar[d]  \\
     S' \ar[r]          & S
 \end{tikzcd}
\]
of \emph{affinoid} adic spaces induce morphisms of adic spaces $\Spa(X',S') \to \Spa(X,S)$.

For the general case, we use the following construction for the definition of a topology.
For all possible commutative squares
\begin{equation} \label{test-square-comp-topology}
 \begin{tikzcd}
     U  \ar[r,open] \ar[d]  & X \ar[d]  \\
     T  \ar[r,open]         & S,
 \end{tikzcd}
\end{equation}
where~$U$ and~$T$ are affinoid and the horizontal maps are open immersions, we consider the associated morphisms $\Spa(U,T) \to \Spa(X,S)$. 
By \cref{lem:Spa-injective} this map is injective, and we can thus treat $\Spa(U,T)$ as a subset of $\Spa(X,S)$.
Now we endow $\Spa(X,S)$ with the final topology for these inclusions.
This means that a subset of $\Spa(X,S)$ is open if and only if its intersection with $\Spa(U,T)$ is open for every square as above.
The first sanity check is the following:

\begin{lem}
    If $X = \Spa(A,A^+)$ and $S = \Spa(R,R^+)$ are affinoid, the above defined topology on $\Spa(X,S)$ coincides with the topology of the adic space $\Spa(A,R^+)$ via the identification from \cref{lem:Spa-barX-affinoid}.
\end{lem}

\begin{proof}
    We consider a subset $\cU \subseteq \Spa(X,S)$.
    A square (\ref{test-square-comp-topology}) comes from a square
    \[
        \begin{tikzcd}
            (B,B^+)         & (A,A^+) \ar[l]    \\
            (D,D^+) \ar[u]  & (R,R^+) \ar[l] \ar[u]
        \end{tikzcd}
    \]
    of Huber pairs and thus induces a morphism of affinoid adic spaces $\Spa(B,D^+) \to \Spa(A,R^+)$ whose underlying map of sets identifies with $\Spa(U,T) \to \Spa(X,S)$.
    Consequently, if~$\cU$ is open in $\Spa(X,S)$ with the topology coming from the structure as an affinoid adic space, then $\cU \cap \Spa(U,T)$ is open.

    Conversely, suppose that $\cU \cap \Spa(U,T)$ is open for all squares~(\ref{test-square-comp-topology}).
    This applies in particular to $\Spa(U,T) = \Spa(X,S)$.
    Hence~$\cU$ is open in $\Spa(X,S)$ with its adic space topology.
\end{proof}

Suppose that the hypotheses of \cref{thm:Solodov} are satisfied for $f \colon X \to S$.
We define a map
\[
    \phi \colon \Spa(X,S) \la \ov X
\]
by sending a root triple $x = (x^\circ,k(x)^+,x^\succ)$ to the unique centre of~$k(x)^+$ in~$\ov X$ lying over~$x^\succ$ (which exists and is unique by the valuative criterion for vertical partial properness \cite{Bauer:valuative-criteria}).
We obtain a commutative diagram
\[
 \begin{tikzcd}
     X  \ar[rr] \ar[dr,"\iota"]  \ar[ddr,"f"']   &           & \ov X \ar[ddl]    \\
     & \Spa(X,S) \ar[d]  \ar[ur,"\phi"] \\
     & S.
 \end{tikzcd}
\]
Here, the map $\iota \colon X \to \Spa(X,S)$ sends~$x$ to $(x^\circ,k(x)^+,f(x))$ (alternatively, identify~$X$ with $\Spa(X,X)$ as in \cref{ex:comp-affinoid} and use the natural map $\Spa(X,X) \to \Spa(X,S)$).
The commutativity of the upper triangle is ensured by the vertical separatedness of $\ov X \to S$.


\begin{prop}[\cite{Solodov-vertical-compactification}, Theorem~4.8] \label{prop:comp-homeo}
    Under the hypotheses of \cref{thm:Solodov}, the above defined map~$\phi$ is a homeomorphism.
\end{prop}

By the proposition the map
$\iota \colon X = \Spa(X,X) \longrightarrow \Spa(X,S)$ is a topological open embedding.
We endow $\Spa(X,S)$ with a structure sheaf by setting 
\[
\cO_{\Spa(X,S)} = \iota_* \cO_X.
\]
The computation of the stalk of $\iota_*\cO_X$ at a point $x \in \Spa(X,S)$ relies on the following lemma.

\begin{lem} \label{lem:stalks-comp}
    Suppose $f \colon X \to S$ satisfies the hypotheses of \cref{thm:Solodov}.
    Let $x \in \Spa(X,S)$ be a point and denote by~$x'$ the minimal vertical generalisation of~$x$ lying in~$X$ (so if $x' \leadsto y \leadsto x$ are vertical specialisations and $x' \ne y$, then~$y$ is not contained in~$X$).
    Then the natural homomorphism $(\iota_*\cO_X)_x \to (\iota_*\cO_X)_{x'} = \cO_{X,x'}$ is an isomorphism.
\end{lem}

\begin{proof}
    We need to show the following:
    For every open neighbourhood $U' \subseteq X$ of~$x'$ there is an open neighbourhood $U \subseteq \Spa(X,S)$ of~$x$ such that $U \cap X \subseteq U'$.

    The subset $\Spa(U',S) \subseteq \Spa(X,S)$ contains~$x$.
    Moreover, by the definition of~$x'$, $x$ is contained in the subset $U \subseteq \Spa(U',S)$ consisting of all points $\bar{z}$ such that for any vertical generalisation $z$ of~$\bar{z}$ we either have $z \in U'$ or $z \notin X$.
    In particular, $U \cap X \subseteq U'$.
    Via the homeomorphism $\Spa(X,S) \cong \ov X$ from \cref{prop:comp-homeo} we can apply \cref{lem:interiors-cover} to $\Spa(X,S)$ telling us that~$U$ is an open neighbourhood of~$x$ in $\Spa(X,S)$.
\end{proof}

\cref{lem:stalks-comp} tells us that $(\Spa(X,S),\iota_*\cO_X)$ is a locally topologically ringed space.
We turn it into a locally $v$-ringed space by defining valuations on the residue fields of the stalks of $\iota_*\cO_X$.
For a point $x = (x^\circ,k(x)^+,x^\succ)$ in $\Spa(X,S)$ the point~$x'$ from \cref{lem:stalks-comp} is of the form $x' = (x^\circ,k(x')^+,x'^\succ)$ where~$k(x')^+$ is the smallest valuation ring satisfying $k(x)^+ \subseteq k(x')^+ \subseteq k(x^\circ)$ and having a centre on~$X$.
By \cref{lem:stalks-comp} we have an equality of residue fields $k(x) = k(x')$ and these are dense in $k(x^\circ)$.
Now we define the valuation ring at~$x$ to be $k(x)^+ \cap k(x)$.

We would like to check that the locally $v$-ringed space $\Spa(X,S)$ is indeed an adic space, i.e.~locally isomorphic to the adic spectrum of a Huber pair.

\begin{thm} \label{thm:Solodov-comp}
    Suppose~$X$ is weakly square complete,~$S$ is square complete and stable, and that the morphism $f \colon X \to S$ is locally of ${}^+$weakly finite type, separated, and taut.
    Then $\Spa(X,S)$ is an adic space isomorphic to the universal vertical compactification~$\overline{X}$ of~$X$ over~$S$.
\end{thm}

\begin{proof}
    We constructed a natural homeomorphism $\phi \colon \Spa(X,S) \to \ov X$ by sending $(x^\circ,k(x)^+,x^\succ)$ to the unique centre of $k(x)^+$ on~$\ov X$ lying over~$x^\succ$ (see \cref{prop:comp-homeo}).
    On the one hand, the universal vertical compactification comes with an open immersion $i \colon X \to \overline{X}$ such that $i_*\cO_X = \cO_{\overline{X}}$.
    On the other hand we defined the structure sheaf on $\Spa(X,S)$ by $\iota_*\cO_X$.
    Via the homeomorphism $\phi$ the two sheaves $i_*\cO_X$ and $\iota_*\cO_X$ are identified.
    It remains to check that the valuations on the residue fields agree.
    For a point $x = (x^\circ,k(x)^+,x^\succ)$ in $\Spa(X,S)$ the valuation ring on $k(x)$ is $k(x)^+ \cap k(x)$ and the same is true for $\phi(x)$.
    This gives both assertions of the theorem.
\end{proof}

\begin{rmk} \label{rmk:hypotheses-compactification}
    The hypotheses of \cref{thm:Solodov-comp} are quite mild.
    Let us discuss the major settings of interest to this paper when they are satisfied.
    First of all, the assumption that~$S$ be sheafy just tells us that the structure presheaf is a sheaf and this continues to hold after base change.
    Tautness is also a minor restriction.
    \begin{enumerate}[(i)]
        \item Suppose that~$S$ and hence also~$X$ are analytic.
              Then~$S$ and~$X$ are automatically square complete.
              Our main application here is to rigid spaces, where $S = \Spa(K)$ is the spectrum of a non-archimedean field and $X \to S$ is locally of finite type.
              These spaces are always stable.
        \item Let us now look at the case of formal schemes, i.e. where $X \to S$ comes from a morphism of formal schemes.
              Also in this situation square completeness is automatic once we require~$X$ to be separated:
              If $x \leadsto v$ is a vertical specialisation and $x \leadsto h$ a horizontal specialisation, we can find an affinoid neighbourhood of~$v$ of the form $\Spa(A)$.
              Then also $x \in \Spa(A)$ and for every intermediate valuation ring $k(x)^+ \subseteq V \subseteq k(x)^\circ$ there is a corresponding horizontal specialisation in $\Spa(A)$.
              By the separatedness assumption, the horizontal specialisations of $x \in X$ are uniquely determined by a valuation ring~$V$ as above, so in particular $h \in \Spa(A)$.
              Since affinoids are square complete, we can find $z \in \Spa(A)$ such that $v \leadsto z$ horizontally and $h \leadsto z$ vertically.

              In \cite{ZavyalovSheafy} Zavyalov shows that strongly rigid-noetherian Huber pairs are sheafy.
              In particular, formal schemes that are locally the formal spectrum of such a ring give rise to stable adic spaces (sending $\Spf(A)$ to $\Spa(A)$).
              This applies to formal schemes that are locally of topologically finite type over a complete microbial valuation ring (e.g. over~$K^\circ$ for a non-archimedean field~$K$).
              So formal models of rigid spaces over a non-archimedean field satisfy the assumptions of \cref{thm:Solodov-comp}.
        \item For discretely ringed adic spaces stability is not an issue.
              Square completeness is satisfied for adic spaces of the form $\Spa(Y,T)$ for a morphism of schemes $Y \to T$.
              Moreover, tautness is certainly satisfied for spaces coming from schemes of finite type over a noetherian base.
              Our main application is to adic spaces of the form $X = \Spa(Y,k)$, where~$k=K^\succ$ is the residue field of a non-archimedean field $(K,K^\circ)$ and $Y$ is a scheme of finite type over~$k$.
              The base~$S$ is usually $S = \Spa(k)$.
    \end{enumerate}
\end{rmk}


We have already noted that for a morphism of \emph{analytic} adic spaces the universal vertical compactification is actually the universal compactification and coincides with Huber's construction in the noetherian setting.
In the opposite corner in the category of adic spaces there are located the discretely ringed adic spaces, and indeed for this type of adic spaces the universal vertical compactification recovers a familiar construction from \cite[\S~3.1]{Temkin2011:RelativeRiemannZariskiSpacesa}:


\begin{ex} \label{example-Spa-schemes}
        Let $f_0\colon X_0\to S_0$ be a morphism of schemes and let 
    \[
        f\colon X=\Spa(X_0, X_0)\la \Spa(S_0, S_0)=S
    \]
    be the associated morphism of discretely ringed adic spaces. Then $\Spa(X, S)$ is isomorphic to the discretely ringed adic space $\Spa(X_0, S_0)$.
    Indeed, by \cref{defi:Spa for schemes}, points of $\Spa(X_0, S_0)$ are triples $(x_0, V, \varphi)$ where $x_0\in X_0$, $V$ is a valuation subring of $k(x_0)$, and $\varphi \colon \Spec(V)\to S_0$ is a map of schemes making the square
\[ 
    \begin{tikzcd}
        \Spec(k(x_0)) \ar[r] \ar[d] & X_0 \ar[d] \\
        \Spec(V)\ar[r,"\varphi"] & S_0
    \end{tikzcd}
\]
commute.
The point $x_0$ lifts to a unique point $x^\circ$ of $X=\Spa(X_0, X_0)$ without vertical generalisations by endowing $k(x_0)$ with the trivial valuation.
Then $k(x^\circ) = k(x^\circ)^+ = k(x_0)$ and we set $k(x)^+ := V$.
Then $\varphi$ induces a map $\Spa(k(x^\circ), V)\to S$.
We denote by $x^\succ \in S$ the image of the closed point.
This determines a root triple $x = (x^\circ,k(x)^+,x^\succ)$, i.e. a point of $\Spa(X,S)$.
The resulting map $\Spa(X_0,S_0) \to \Spa(X,S)$ is bijective and the following observation shows that it respects the adic space structure by reducing to the affine case:
For every test square
\[
    \begin{tikzcd}
        U_0 \ar[r]  \ar[d]  & X_0   \ar[d] \\
        T_0 \ar[r]          & S_0
    \end{tikzcd}
\]
for affine schemes~$U_0$ and~$T_0$ with open immersions $U_0 \to X_0$ and $T_0 \to S_0$, the induced morphism of adic spaces $\Spa(U_0,T_0) \to \Spa(X_0,S_0)$ is an open immersion.
\end{ex}


\subsection{Functoriality of the vertical compactification}

We end this section by exploring the functoriality of the universal vertical compactification.

\begin{prop} \label{prop:Spa-functoriality}
Suppose that we are given a commutative square
\[
 \begin{tikzcd}
     Y     \ar[r,"f_X"]   \ar[d]  & X   \ar[d]  \\
     T      \ar[r,"f_S"]          & S
 \end{tikzcd}
\]
of adic spaces such that $f_S$ is adic and the vertical morphisms satisfy the assumptions of \cref{thm:Solodov-comp}.
Then there is a unique morphism of adic spaces $f \colon \Spa(Y,T) \to \Spa(X,S)$ making, as the dashed arrow, the diagram below commute.
\[
 \begin{tikzcd}[column sep= 1cm]
     Y          \ar[r,"f_X"]   \ar[d]  & X   \ar[d]  \\
     \Spa(Y,T)  \ar[r,dashed,"f"] \ar[d] & \Spa(X,S) \ar[d] \\
     T          \ar[r,"f_S"]          & S
 \end{tikzcd}
\]
In terms of root triples,~$f$ maps a point $y = (y^\circ,k(y)^+,y^\succ)$ to $(f_X(y^\circ)^\circ,k(y)^+ \cap k(f_X(y^\circ)^\circ),f_S(y^\succ))$ (i.e. the underlying map of sets is the one described in \cref{sec:root-triples}).
\end{prop}

\begin{proof}
    We consider the solid arrow diagram
    \[
    \begin{tikzcd}
    & Y \ar[dl] \ar[d] \ar[r,"f_X"] & X \ar[d] \\
    \Spa(Y,T) \ar[r,dashed] \ar[dr] & \Spa(X,S) \times_S T \ar[r,"\pr"] \ar[d,"\rm vpp"] & \Spa(X,S) \ar[d,"\rm vpp"] \\
    & T \ar[r,"f_S"] & S.
    \end{tikzcd}
    \]
    The fibre product exists and $\Spa(X,S) \times_S T \to T$ is vertically partially proper by our assumption that $T \to S$ is adic.
    There is a unique dashed arrow by the universal property of the universal vertical compactification~$\Spa(Y,T)$.
    Composing it with the projection $\pr$,
    we obtain the desired morphism.
    
    By the valuative criterion for vertical partial properness \cite{Bauer:valuative-criteria}, the dotted map sends $y \in \Spa(Y,T)$ to the unique centre of $k(y)^+$ on $\Spa(X,S) \times_S T$ lying over~$T$.
    The projection to $\Spa(X,S)$ then maps it to the unique centre of $k(y)^+$ on $\Spa(X,S)$ lying over $f_S(y^\succ)$.
    This centre identifies with the root triple $(f_X(y^\circ)^\circ,k(y)^+ \cap k(f_X(y^\circ)^\circ),f_S(y^\succ))$.
\end{proof}

\begin{prop} \label{prop:Spa-iso-vpp}
    Suppose that the hypotheses of \cref{thm:Solodov} are satisfied for $X \to T$ and $T \to S$ is vertically partially proper.
    Then $\Spa(X,T) \to T \to S$ is the universal vertical compactification of~$X$ over~$S$.
\end{prop}

\begin{proof}
    We first note that $\Spa(X,T) \to S$ is vertically partially proper, so it is a vertical compactification of~$X$ over~$S$.
    In order to check universality, we consider a solid arrow commutative diagram
    \[
        \begin{tikzcd}
             X  \ar[drr,bend left]   \ar[dr,dashed] \ar[d]  \\
             \Spa(X,T) \ar[r,dashed] \ar[dr] & Y_T \ar[r] \ar[d,"\rm vpp"] \ar[r] & Y \ar[d,"\rm vpp"] \\
             & T \ar[r] & S,
        \end{tikzcd}
    \]
    where~$Y \to S$ is a vertically partially proper morphism and~$Y_T = Y \times_S T$.
    The universal property of the fibre product gives the diagonal dashed arrow and then the universal property for $\Spa(X,T)$ gives the horizontal dashed arrow.
    Composing with $Y_T \to Y$ we obtain the unique map required in the universal property of the universal vertical compactification of~$X$ over~$S$.
\end{proof}

What we have seen so far in this subsection relies entirely on the universal property of the universal vertical compactification.
We did not have to know anything about the construction or further geometric properties.
This will change with the following proposition where we will actually argue using the geometry of the spaces involved.

\begin{prop} \label{prop:comp-fibre-product}
    We consider a commutative diagram
        \[ 
            \begin{tikzcd}
                Y'\ar[r] \ar[d] & Y \ar[d] \\
                X'\ar[r] \ar[d] & X \ar[d] \\
                S'\ar[r]  & S
            \end{tikzcd}
        \]
    of adic spaces with $Y' \simeq X' \times_X Y$.
    Assuming that all vertical morphisms to~$S$ and~$S'$ satisfy the hypotheses of \cref{thm:Solodov-comp} and the horizontal morphisms  are adic, it follows that the natural morphism
    \[
     \Spa(Y',S') \longrightarrow \Spa(X',S') \times_{\Spa(X,S)} \Spa(Y,S)
    \]
    is an isomorphism of adic spaces.
\end{prop}

\begin{proof}
    The fibre product $\Spa(X',S') \times_{\Spa(X,S)} \Spa(Y,S)$ exists because $\Spa(Y,S) \to \Spa(X,S)$ is locally of ${}^+$weakly finite type (being vertically partially proper) and $\Spa(X',S') \to \Spa(X,S)$ is adic because $X' \to X$ is adic.
    By functoriality of $\Spa(-,-)$ as in \cref{prop:Spa-functoriality}, we obtain the natural morphism
    \[
       \varphi \colon \Spa(Y',S') \longrightarrow \Spa(X',S') \times_{\Spa(X,S)} \Spa(Y,S).
    \]
    Using the explicit description via root triples, we check that~$\varphi$ is bijective.
    Now, it suffices to show that $\varphi$ is locally an open immersion.

    We first reduce to the case that $S$ and $S'$ are affinoid as follows. 
    In the situation of \cref{thm:Solodov-comp}, the hypotheses of \cref{prop:check-univprop-comp} 
    are satisfied (see the proof of \cref{prop:check-univprop-comp} in \cite[Theorem~3.29]{Solodov-vertical-compactification}).
    This implies that not only is $\Spa(X,S)$ the universal compactification of~$X$ over~$S$, but also for every open subspace $U \subseteq S$ the base change $\Spa(X,S) \times_S U$ is the universal compactification of $X \times_S U$ over~$U$.
    Applying this reasoning to $X \to S$, $X' \to S'$, $Y \to S$, and $Y' \to S'$, we may now assume that $S$ and $S'$ are affinoid. 

    We choose compatible affinoid open covers $\{U_i\}_{i \in I}$, $\{U'_j\}_{j \in J}$, and $\{V_k\}_{k \in K}$ of~$X$,~$X'$, and~$Y$, respectively, such that for all $i \in I$ we have
    \[
        U_i \times_X X' = \bigcup_{j \to i} U'_j, \qquad U_i \times_X Y = \bigcup_{k \to i} V_k.
    \]
    We can apply \cref{lem:compatible-covers} below to the squares
    \[
        \begin{tikzcd}
            X' \ar[r] \ar[d] & X \ar[d] \\
            \Spa(X',S')  \ar[r] & \Spa(X,S)
        \end{tikzcd} \qquad
             \begin{tikzcd}
            Y \ar[r] \ar[d] & X \ar[d] \\
            \Spa(Y,S)  \ar[r] & \Spa(X,S)
        \end{tikzcd}
    \]
    and the open subsets $\ov U_i^\circ \subseteq \Spa(X,S)$ from \cref{lem:interiors-cover}, noting that $\ov U_i^\circ \cap X = U_i$.
    Consequently, for fixed~$i$, the open subsets 
    \begin{itemize}
        \item 
        $\widetilde{U}'_j := (\ov U^\circ_i \times_{\Spa(X,S)} \Spa(X',S')) \cap \ov U'^\circ_j$ for $j \to i$ and
        \item 
        $\widetilde{V}_k := (\ov U^\circ_i \times_{\Spa(X,S)} \Spa(Y,S)) \cap \ov V_k^\circ$ for $k \to i$
    \end{itemize}  
    cover $\ov U^\circ_i \times_{\Spa(X,S)} \Spa(X',S')$ and $\ov U^\circ_i \times_{\Spa(X,S)} \Spa(Y,S)$, respectively.
    Moreover $\widetilde{U}'_j \cap X' = U'_j$ and $\widetilde{V}_k \cap Y = V_k$ and for $j \to i$ and $k \to i$, the images of~$\widetilde{U}'_j$ and $\widetilde{V}_k$ are contained in~$\ov U^\circ_i$.
    Hence, the fibre products $\widetilde{U}'_j \times_{\ov U^\circ_i} \widetilde{V}_k$ cover $\Spa(X',S') \times_{\Spa(X,S)} \Spa(Y,S)$ and we can apply \cref{lem:compatible-covers} again, this time to the square
    \[
        \begin{tikzcd}
            Y' \ar[r,"\id"] \ar[d] & X' \times_X Y \ar[d] \\
            \Spa(Y',S') \ar[r,"\varphi"] & \Spa(X',S') \times_{\Spa(X,S)} \Spa(Y,S)
        \end{tikzcd}
    \]
    with the open subsets $\widetilde{U}'_j \times_{\ov U^\circ_i} \widetilde{V}_k$, and the trivial open cover of the affinoid $V'_{j,k} = U'_j \times_{U_i} V_k$ by itself (for all $(j,k) \in J \times_I K$).
    In this way we construct an open cover of $\Spa(Y',S')$ by subsets $\widetilde{V}'_{j,k}$  lying in $\Spa(V'_{j,k},S')$ and fitting into the diagram
        \[
        \begin{tikzcd}
            & \Spa(V'_{j,k},S') \arrow[rr,"\varphi_{j,k}"] \arrow[dd]
            && \Spa(U'_j,S') \times_{\Spa(U_i,S)} \Spa(V_k,S) \arrow[dd]
            \\
            \widetilde{V}'_{j,k} \arrow[rr,dashed, crossing over]  \arrow[ur,open] \arrow[dr,open]
            &&  \widetilde{U}'_j \times_{\ov{U}_i^\circ} \widetilde{V}_k \arrow[ur,open] \arrow[dr,open]
            \\
            & \Spa(Y',S') \arrow[rr,"\varphi"] 
            && \Spa(X',S') \times_{\Spa(X,S)} \Spa(Y,S) 
        \end{tikzcd}   
    \]
    where all diagonal maps are open immersions. The square with the morphisms $\ph$ and $\varphi_{j,k}$ commutes by the functoriality of $\Spa(-,-)$, hence both morphisms agree on $\widetilde{V}'_{j,k}$.
    Since the subsets~$\widetilde{V}'_{j,k}$ cover $\Spa(Y',S')$, if we knew that~$\varphi_{j,k}$ is an isomorphism, we could conclude that~$\varphi$ is a bijective open immersion, hence an isomorphism.
    
    So finally, we need to treat the case where $S = \Spa(R,R^+)$, $S' = \Spa(R',R'^+)$, $X = \Spa(A,A^+)$, $Y = \Spa(B,B^+)$, and $X' = \Spa(A',A'^+)$ are affinoid.
    Then
    \begin{align*}
        \Spa(X',S') \times_{\Spa(X,S)} \Spa(Y,S) &=  \Spa(A',R'^+) \times_{\Spa(A,R^+)} \Spa(B,R^+) \\
        & = \Spa(A' \otimes_A B,R'^+) 
        = \Spa(Y',S'). \qedhere
    \end{align*}    
\end{proof}

\begin{lem} \label{lem:compatible-covers}
    Let $T$ be an affinoid adic space. Let $Z \to T$ be a morphism of adic spaces satisfying the hypotheses of \cref{thm:Solodov-comp}.
    Suppose we are given a diagram
    \[
        \begin{tikzcd}
            Z   \ar[r,"\varphi_0"]  \ar[d,open] & W_0 \ar[d,open] \\
            \Spa(Z,T) \ar[r,"\varphi"]                  & W
        \end{tikzcd}
    \]
    of adic spaces, with an open immersion $W_0 \to W$ and an open subspace $U \subseteq W$. 
    If $\varphi_0^{-1}(U \cap W_0) = \bigcup_{j \in J} V_j$ is a covering by open affinoid subspaces $V_j \subseteq Z$, then, with the open $\ov V_j^\circ \subseteq \Spa(Z,T)$ as in \Cref{lem:interiors-cover} we have 
    \[
        \varphi^{-1}(U) \subseteq \bigcup_{j \in J} \ov V_j^\circ.
    \]    
\end{lem}
\begin{proof}
    For~$z \in \varphi^{-1}(U)$ let $z_0 \leadsto z$ be the minimal vertical generalisation of~$z$ contained in~$Z$.
    Then $\varphi(z_0) = \varphi_0(z_0) \in U \cap W_0$.
    We choose $j \in J$ with $z_0 \in V_j$.
    Then~$z$ is contained in~$\ov V_j^\circ$.
\end{proof}

\begin{cor} \label{cor:functoriality-compactification}
    Let $X \to S$ be as in \cref{thm:Solodov-comp}.
    For any adic morphism $S' \to S$ with $S'$ stable and square complete, the base change $\Spa(X,S)\times_S S'$ is the universal vertical compactification of $X \times_S S'$ over~$S'$.
    In other words, $\Spa(X\times_S S',S')$ is an adic space with a natural isomorphism
    \[
     \Spa(X \times_S S',S') \overset{\sim}{\longrightarrow} \Spa(X,S)\times_S S'.
    \]
\end{cor}


\begin{proof}
We want to apply \cref{prop:comp-fibre-product} to the diagram
\[
    \begin{tikzcd}
        X \times_S S'\ar[r] \ar[d] & X\ar[d] \\
        S'\ar[r] \ar[d] & S \ar[d] \\
        S'\ar[r] & S.
    \end{tikzcd}
\]
For doing so we still need to check that $X \times_S S' \to S'$ satisfies the hypotheses of \cref{thm:Solodov-comp}, i.e. that $X \times_S S'$ is weakly square complete.
This follows from the subsequent \cref{lem:square-complete-stable}.
\end{proof}

\begin{lem} \label{lem:square-complete-stable}
    Let
    \[
        \begin{tikzcd}
            X'  \ar[r]  \ar[d]  & X \ar[d]  \\
            S'  \ar[r]          & S
        \end{tikzcd}
    \]
    be a cartesian diagram of adic spaces, where~$S' \to S$ is adic, $S$ is square complete, and~$X$ and~$S'$ are weakly square complete.
    Then~$X'$ is weakly square complete.
\end{lem}

\begin{proof}
    We consider a diagram
         \[
      \begin{tikzcd}
       x'_v									\\
       x'	\ar[r,rightsquigarrow]	\ar[u,rightsquigarrow]	& x'_h	
      \end{tikzcd}
     \]
     in~$X'$, where $x' \leadsto x'_v$ is a vertical and $x' \leadsto x'_h$ a horizontal specialisation.
     We need to find $x'_{vh} \in X'$, a horizontal specialisation $x'_v \leadsto x'_{vh}$, and a vertical specialisation $x'_h \leadsto x'_{vh}$ completing the square.
     This is possible once we find a common affinoid neighbourhood   of~$x'$, $x'_h$, and~$x'_v$, as affinoids are square complete.
     
     Since~$X$ and~$S'$ are weakly square complete, the image diagrams in~$X$ and~$S'$ can be completed to squares
          \[
      \begin{tikzcd}
       x_v	\ar[r,rightsquigarrow]							& x_{vh}			\\
       x	\ar[r,rightsquigarrow]	\ar[u,rightsquigarrow]	& x_h	\ar[u,rightsquigarrow]
      \end{tikzcd} \qquad
      \begin{tikzcd}
       s'_v	\ar[r,rightsquigarrow]							& s'_{vh}			\\
       s'	\ar[r,rightsquigarrow]	\ar[u,rightsquigarrow]	& s'_h	\ar[u,rightsquigarrow]
      \end{tikzcd}
     \]
     with vertical and horizontal specialisations according to the definition of weak square completeness.
     Since~$S$ is even square complete, the images of both squares in~$S$ are the same.
     Choosing small enough affinoid neighbourhoods $U$, $V'$, $V$ of~$x_{vh}$, $s'_{vh}$, and their image $s_{vh}$ in~$S$, respectively, the fibre product $U' = U \times_{V} V'$ is an affinoid open in $X'$ containing $x'$, $x'_v$ and $x'_h$.
\end{proof}

\begin{rmk}
    In \cref{prop:comp-fibre-product} and \cref{cor:functoriality-compactification} the assumption that the relevant morphisms be adic is not only needed to ensure that the fibre product exists.
    For instance, the need for $S' \to S$ to be adic in \cref{cor:functoriality-compactification} is illustrated by the following example.
    Let $S'$ be the adic spectrum of a non-archimedean field $(K,K^\circ)$.
    We denote by~$S$ the spectrum of the same valued field $(K,K^\circ)$ but equipped with the discrete topology.
    The resulting morphism $S' \to S$ is not adic.
    For $X = \Spa(K,K)$ (with discrete topology on~$K$) the universal vertical compactification of~$X$ over~$S$ equals
    \[
     \Spa(X,S) = \Spa(K,K^\circ) = S.
    \]
    Its base change to~$S'$ equals~$S'$.
    But $X' = X \times_S S'$ is empty, so $\Spa(X,S) \times_S S'$ cannot be the universal vertical compactification of~$X'$ over~$S'$.
\end{rmk}


\begin{cor} \label{cor:base-change-comp}
    Let~$K$ be a non-archimedean field with residue field~$k$ and~$\fX$ a taut and separated formal scheme locally of finite type over~$K^\circ$.
    Let $\fX_\rig$ be its rigid generic fibre and~$\fX_k$ its special fibre.
    Then
    \begin{align*}
     \Spa(\fX^{\rm ad},K^\circ)_\eta &= \Spa(\fX_\rig,K),   \\
     \Spa(\fX^{\rm ad},K^\circ)_k &= \Spa(\fX_k,k).
    \end{align*}
\end{cor}

\begin{proof}
    This follows by applying \cref{cor:functoriality-compactification} to the base change of $\fX^{\rm ad} \to \Spa(K^\circ)$ via $\Spa(K) \to \Spa(K^\circ)$ and $\Spa(k) \to \Spa(K^\circ)$, respectively.
\end{proof}

\begin{ex} \label{ex:comp-disk}
    We illustrate the identifications from  \cref{cor:base-change-comp} by going through the constructions when
    \[
        \fX = \bA^1_{K^\circ} =  \Spf (K^\circ \langle T \rangle)
    \]
    is the formal affine line.
    Its generic fibre is the unit disc
    \[
        \fX_\eta = \bD_K = \Spa (K \langle T \rangle),
    \]
    whose compactification $\Spa(\bD_K,K)$ has one additional point~$x_{1^+}$ corresponding to the valuation of $K \langle T \rangle$ sending~$T$ to a value that is infinitesimally bigger than~$1$.
    It can be written as the composition of the Gauß valuation (of radius~$1$) with the valuation on its specialisation field $k(T)$ corresponding to $\infty \in \bP^1_k$.
    The special fibre is the affine line
    \[
        \fX_k = \bA^1_k = \Spec (k[T])
    \]
    with universal vertical compactification $\Spa(k[T],k)$.
    Also here we get one additional point~$x_{\infty}$, namely the valuation of $k(T)$ corresponding to the point $\infty \in \bP^1_k$.
    Now according to \cref{cor:base-change-comp}, the universal vertical compactification $\Spa(\bA^1_{K^\circ},K^\circ)$ has the two additional points~$x_{1^+}$ and~$x_\infty$.
    Note that~$x_\infty$ is a horizontal specialisation of~$x_{1^+}$ and if we were to extend the specialisation map from \cref{ss:adic-spaces} to universal vertical compactifications, we would send~$x_{1^+}$ to~$x_\infty$.
\end{ex}


